% Options for packages loaded elsewhere
\PassOptionsToPackage{unicode}{hyperref}
\PassOptionsToPackage{hyphens}{url}
\documentclass[
]{article}
\usepackage{xcolor}
\usepackage{amsmath,amssymb}
\setcounter{secnumdepth}{-\maxdimen} % remove section numbering
\usepackage{iftex}
\ifPDFTeX
  \usepackage[T1]{fontenc}
  \usepackage[utf8]{inputenc}
  \usepackage{textcomp} % provide euro and other symbols
\else % if luatex or xetex
  \usepackage{unicode-math} % this also loads fontspec
  \defaultfontfeatures{Scale=MatchLowercase}
  \defaultfontfeatures[\rmfamily]{Ligatures=TeX,Scale=1}
\fi
\usepackage{lmodern}
\ifPDFTeX\else
  % xetex/luatex font selection
\fi
% Use upquote if available, for straight quotes in verbatim environments
\IfFileExists{upquote.sty}{\usepackage{upquote}}{}
\IfFileExists{microtype.sty}{% use microtype if available
  \usepackage[]{microtype}
  \UseMicrotypeSet[protrusion]{basicmath} % disable protrusion for tt fonts
}{}
\makeatletter
\@ifundefined{KOMAClassName}{% if non-KOMA class
  \IfFileExists{parskip.sty}{%
    \usepackage{parskip}
  }{% else
    \setlength{\parindent}{0pt}
    \setlength{\parskip}{6pt plus 2pt minus 1pt}}
}{% if KOMA class
  \KOMAoptions{parskip=half}}
\makeatother
\usepackage{color}
\usepackage{fancyvrb}
\newcommand{\VerbBar}{|}
\newcommand{\VERB}{\Verb[commandchars=\\\{\}]}
\DefineVerbatimEnvironment{Highlighting}{Verbatim}{commandchars=\\\{\}}
% Add ',fontsize=\small' for more characters per line
\newenvironment{Shaded}{}{}
\newcommand{\AlertTok}[1]{\textcolor[rgb]{1.00,0.00,0.00}{\textbf{#1}}}
\newcommand{\AnnotationTok}[1]{\textcolor[rgb]{0.38,0.63,0.69}{\textbf{\textit{#1}}}}
\newcommand{\AttributeTok}[1]{\textcolor[rgb]{0.49,0.56,0.16}{#1}}
\newcommand{\BaseNTok}[1]{\textcolor[rgb]{0.25,0.63,0.44}{#1}}
\newcommand{\BuiltInTok}[1]{\textcolor[rgb]{0.00,0.50,0.00}{#1}}
\newcommand{\CharTok}[1]{\textcolor[rgb]{0.25,0.44,0.63}{#1}}
\newcommand{\CommentTok}[1]{\textcolor[rgb]{0.38,0.63,0.69}{\textit{#1}}}
\newcommand{\CommentVarTok}[1]{\textcolor[rgb]{0.38,0.63,0.69}{\textbf{\textit{#1}}}}
\newcommand{\ConstantTok}[1]{\textcolor[rgb]{0.53,0.00,0.00}{#1}}
\newcommand{\ControlFlowTok}[1]{\textcolor[rgb]{0.00,0.44,0.13}{\textbf{#1}}}
\newcommand{\DataTypeTok}[1]{\textcolor[rgb]{0.56,0.13,0.00}{#1}}
\newcommand{\DecValTok}[1]{\textcolor[rgb]{0.25,0.63,0.44}{#1}}
\newcommand{\DocumentationTok}[1]{\textcolor[rgb]{0.73,0.13,0.13}{\textit{#1}}}
\newcommand{\ErrorTok}[1]{\textcolor[rgb]{1.00,0.00,0.00}{\textbf{#1}}}
\newcommand{\ExtensionTok}[1]{#1}
\newcommand{\FloatTok}[1]{\textcolor[rgb]{0.25,0.63,0.44}{#1}}
\newcommand{\FunctionTok}[1]{\textcolor[rgb]{0.02,0.16,0.49}{#1}}
\newcommand{\ImportTok}[1]{\textcolor[rgb]{0.00,0.50,0.00}{\textbf{#1}}}
\newcommand{\InformationTok}[1]{\textcolor[rgb]{0.38,0.63,0.69}{\textbf{\textit{#1}}}}
\newcommand{\KeywordTok}[1]{\textcolor[rgb]{0.00,0.44,0.13}{\textbf{#1}}}
\newcommand{\NormalTok}[1]{#1}
\newcommand{\OperatorTok}[1]{\textcolor[rgb]{0.40,0.40,0.40}{#1}}
\newcommand{\OtherTok}[1]{\textcolor[rgb]{0.00,0.44,0.13}{#1}}
\newcommand{\PreprocessorTok}[1]{\textcolor[rgb]{0.74,0.48,0.00}{#1}}
\newcommand{\RegionMarkerTok}[1]{#1}
\newcommand{\SpecialCharTok}[1]{\textcolor[rgb]{0.25,0.44,0.63}{#1}}
\newcommand{\SpecialStringTok}[1]{\textcolor[rgb]{0.73,0.40,0.53}{#1}}
\newcommand{\StringTok}[1]{\textcolor[rgb]{0.25,0.44,0.63}{#1}}
\newcommand{\VariableTok}[1]{\textcolor[rgb]{0.10,0.09,0.49}{#1}}
\newcommand{\VerbatimStringTok}[1]{\textcolor[rgb]{0.25,0.44,0.63}{#1}}
\newcommand{\WarningTok}[1]{\textcolor[rgb]{0.38,0.63,0.69}{\textbf{\textit{#1}}}}
\usepackage{longtable,booktabs,array}
\newcounter{none} % for unnumbered tables
\usepackage{calc} % for calculating minipage widths
% Correct order of tables after \paragraph or \subparagraph
\usepackage{etoolbox}
\makeatletter
\patchcmd\longtable{\par}{\if@noskipsec\mbox{}\fi\par}{}{}
\makeatother
% Allow footnotes in longtable head/foot
\IfFileExists{footnotehyper.sty}{\usepackage{footnotehyper}}{\usepackage{footnote}}
\makesavenoteenv{longtable}
\setlength{\emergencystretch}{3em} % prevent overfull lines
\providecommand{\tightlist}{%
  \setlength{\itemsep}{0pt}\setlength{\parskip}{0pt}}
\usepackage{bookmark}
\IfFileExists{xurl.sty}{\usepackage{xurl}}{} % add URL line breaks if available
\urlstyle{same}
\hypersetup{
  hidelinks,
  pdfcreator={LaTeX via pandoc}}

\author{}
\date{}

\begin{document}

\section{The Lean Mathematical
Library}\label{the-lean-mathematical-library}

The mathlib Community*

\section{Abstract}\label{abstract}

This paper describes mathlib, a community-driven effort to build a
unified library of mathematics formalized in the Lean proof assistant.
Among proof assistant libraries, it is distinguished by its dependently
typed foundations, focus on classical mathematics, extensive hierarchy
of structures, use of large- and small-scale automation, and distributed
organization. We explain the architecture and design decisions of the
library and the social organization that has led to its development.

CCS Concepts • Mathematics of computing → Mathematical software; •
Security and privacy → Logic and verification.

Keywords Lean, mathlib, formal library, formal proof

\section{ACM Reference Format:}\label{acm-reference-format}

The mathlib Community. 2020. The Lean Mathematical Library. In
Proceedings of the 9th ACM SIGPLAN International Conference on Certified
Programs and Proofs (CPP '20), January 20--21, 2020, New Orleans, LA,
USA. ACM, New York, NY, USA, 15 pages.
https://doi.org/10.1145/3372885.3373824

\section{1 Introduction}\label{introduction}

Since the first mechanized proof-checking systems were released, there
has been continual effort to develop libraries of formal definitions and
proofs. Every major system has at least one library that can serve as a
base for further formalizations. The contents, organizations, and
purposes of these libraries vary considerably. Some are small, static
cores bundled with the system itself; some are sprawling repositories
that solicit contributions like a journal; some focus on a narrow
subject area, while others are more broad-minded in their topics.

*This paper describes a community effort with 73 contributors, listed in
the acknowledgments. Rather than attributing authorship to a subset, we
have decided to release this paper under a collective pseudonym.

Permission to make digital or hard copies of all or part of this work
for personal or classroom use is granted without fee provided that
copies are not made or distributed for profit or commercial advantage
and that copies bear this notice and the full citation on the first
page. Copyrights for components of this work owned by others than the
author(s) must be honored. Abstracting with credit is permitted. To copy
otherwise, or republish, to post on servers or to redistribute to lists,
requires prior specific permission and/or a fee. Request permissions
from permissions@acm.org. CPP '20, January 20--21, 2020, New Orleans,
LA, USA

© 2020 Copyright held by the owner/author(s). Publication rights
licensed to ACM.

ACM ISBN 978-1-4503-7097-4/20/01\ldots\$15.00

https://doi.org/10.1145/3372885.3373824

This paper describes mathlib, a formal library developed for the Lean
proof assistant {[}20{]}. As a community-driven effort with dozens of
contributors, there is no central organization to mathlib; it has arisen
from the desires of its users to develop a repository of formal
mathematical proofs. We are certainly not the first to profess this goal
{[}1{]}, nor is our library particularly large in comparison to others.
However, its organizational structure, focus on classical mathematics,
and inclusion of automation distinguish it in the space of proof
assistant libraries. We aim here to explain our design decisions and the
ways in which mathlib has been put to use.

In contrast to most modern proof assistant libraries, many of the
contributors to mathlib have an academic background in pure mathematics.
This has significantly influenced the contents and direction of the
library. It is a goal of many in the community to support the
formalization of modern, research-level mathematics, and various
projects discussed in Section 7.2 suggest that we are approaching this
point.

\section{1.1 A History of mathlib and Lean
3}\label{a-history-of-mathlib-and-lean-3}

The Lean project was started by Leonardo de Moura in 2013 {[}20{]}. Its
most recent version, Lean 3, was released in early 2017 {[}22{]}. A new
version is under development {[}57{]}. In the summer of 2017, much of
the core library for Lean was factored out of the system repository.
This code became the base for mathlib. The project was initially led by
Mario Carneiro and Johannes Hölzl at Carnegie Mellon University, and
attracted a growing number of users, who were drawn by the system
documentation {[}5{]}, the library's focus on classical mathematics, and
the real-time chat on Zulip.

Over the two years since the split, mathlib has grown from 15k to 140k
lines of code (LOC), excluding blank lines and comments. Contributions
have been made by 73 people and are managed by a team of 11 maintainers.
It is the de facto standard library for both programming and proving in
Lean 3. The surrounding community has developed infrastructure,
promotional materials, and university courses based on mathlib.

\section{2 Lean}\label{lean}

Like Coq, Lean uses a system of dependent types based on the calculus of
inductive constructions (CIC) {[}52{]}. Lean has a noncumulative
hierarchy of universes Prop, Type, Type 1, Type 2, Type 3 \ldots{} The
bottom universe Prop is impredicative, meaning that quantification of a
Prop over a larger type is a Prop, and proof-irrelevant, meaning that
all proofs of the same proposition are definitionally (or judgmentally)
equal.

Lean adds two axioms to the type theory. The first axiom, classical.
choice, produces a term of type T from a (propositional) proof that T is
nonempty. This is a type-theoretic statement that implies the
set-theoretic axiom of choice. The second axiom, propext, states that a
bi-implication between two propositions implies that these propositions
are equal.

Lean also extends the CIC with quotient types. Given an equivalence
relation on a type, the quotient type identifies related terms. This
feature allows us to work with equality of objects in the quotient and
removes the need for setoids.

A key feature of Lean 3 is its metaprogramming framework {[}22{]}. This
allows users to write tactics, commands, and other tools for
manipulating the Lean environment directly in the language of Lean
itself. In mathlib, we make extensive use of metaprogramming both for
powerful automation and specialized commands to ease various tasks when
formalizing. We will discuss this more in Section 6.

\section{3 Contents of mathlib}\label{contents-of-mathlib}

The mathlib library is designed as a basis for research level
mathematics, as well as a standard library for programming in Lean. We
build mathlib on top of a small core library, shipped with Lean, which
contains 19k LOC. The core library sets up the metaprogramming and
tactic framework for Lean, but it also contains much of the algebraic
hierarchy and the definitions of basic datatypes, like N, Z, and list.

Currently, much of mathlib consists of undergraduate level mathematics.
Table 1 gives the size of all top-level directories in the src
directory. This gives a rough idea of the relative sizes of the
different topics in mathlib, but the contents of some directories
require some extra explanation.

In the algebra library, we expand on the core library algebraic
hierarchy, from semigroup to linear\_ordered\_field, module, and more
(Section 4.1). The results in this folder are focused on computing with
elements in these algebraic structures. The structural theory of these
algebraic objects is developed in the folders group\_theory,
ring\_theory, field\_theory, and linear\_algebra.

The data folder contains the definitions and properties of data
structures, including the number systems \(\mathbb{Q}\) , \(\mathbb{R}\)
, and \(\mathbb{C}\) , sets and subtypes, partial and finitely supported
functions, polynomials, lists, multisets, and vectors.

The folders meta and tactic use Lean's metaprogramming framework
(Section 6) to define custom tactics. The category folder defines
infrastructure for categorical programming, as used in Haskell; this is
not to be confused with the category\_theory folder, which develops the
mathematical theory.

Quotients are heavily used in mathlib to avoid the need for setoids.
They are used to define multisets as lists up to permutation, which are
in turn used to define finite sets as multisets without duplicates.
Quotients are also frequently used in algebra, for example for the
definition of a quotient group, the tensor product of modules or the
colimit of rings. We also use quotients when defining the Stone-Čech
compactification and Cauchy completion, the latter of which is used to
define the real and \(p\) -adic numbers. Quotients are also used to
define cardinals (as types modulo equivalence) and ordinals (as
well-ordered types modulo order-isomorphism).

Subdirectory

LOC

Declarations

data

41849

10695

topology

17382

2709

tactic

12184

1679

algebra

9830

2794

analysis

7962

1237

order

6526

1542

category\_theory

6299

1560

set\_theory

6163

1394

measure\_theory

6113

926

ring\_theory

5683

1080

linear\_algebra

4511

805

computability

4205

575

group\_theory

4191

1094

category

1770

389

number\_theory

1394

228

logic

1195

403

field\_theory

1002

121

geometry

848

70

meta

784

135

algebraic\_geometry

194

29

140085

34168

Table 1. Lines of code, excluding white space and comments, in the
top-level directories of the mathlib source code as of December 12,
2019.

The extensive topology library includes theories about uniform spaces,
metric spaces, and algebraic topological spaces such as topological
groups and rings. A more novel feature is the definition of the
Gromov--Hausdorff space, i.e.~the space of all nonempty compact metric
spaces up to isometry, with its natural (Gromov--Hausdorff) distance,
and the proof that it is a Polish space. \(^{1}\)

The definition of a manifold in mathlib is very general compared to
other known formalizations {[}38, 51{]}. In particular, the base field
for differentiable manifolds can be any non-discrete normed field,
including \(\mathbb{R}\) , \(\mathbb{C}\) , and \(\mathbb{Q}_p\) .
Arbitrary models and structure groupoids can be used. Examples currently
include \(C^k\) and \(C^\infty\) manifolds, possibly with boundary or
corners. Potential future examples include analytic manifolds, contact
or symplectic manifolds, and translation surfaces.

The category theory library includes (co)limits, monadic adjunctions,
and monoidal categories. It uses extensive automation to hide easy
proofs (Section 6.2). The theory established here is used in other parts
of the library, for example describing the Giry monad in measure theory,
or showing that products and equalizers in complete separated uniform
spaces can be calculated in the underlying uniform spaces.

Since Lean's type theory contains a universe hierarchy, it has a higher
consistency strength than ZFC, and the set theory library defines a
model of ZFC. Additionally, cardinals and ordinals are defined using
quotients of a universe, along with related notions like ordinal
arithmetic, \(\aleph_{\alpha}\) , cofinality, inaccessible cardinals,
etc. This machinery is used in other parts of the library (Section 5.3).

Structural results about groups, rings, and fields include Sylow's
theorems, the law of quadratic reciprocity, integral and perfect
closures, and Hilbert's basis theorem. The linear algebra library is
described in detail in Section 5.

A library on computability theory {[}16{]} proves the undecidability of
the halting problem, and a library on the \(p\) -adic numbers proves
Hensel's lemma {[}44{]}.

Analysis is a weaker point of mathlib, although the Fréchet derivative
and the Bochner integral have been formalized, as well as basic
properties of trigonometric functions.

\section{4 Type Classes}\label{type-classes}

Type classes are predicates or extra data attached to types, which are
systematically inferred by a Prolog-like backtracking search. The idea
of using type classes to permit polymorphism over types with a
particular structure was pioneered in Haskell {[}59{]}, and the
usefulness of type classes for organizing mathematical structures and
theorems on these structures was recognized in Isabelle {[}31, 60{]} and
Coq {[}56{]}. Lean's core library builds on these ideas by using type
classes ubiquitously, including for:

\begin{itemize}
\tightlist
\item
  notations, e.g.~the has\_add \(\alpha\) type class that is inferred
  when the + symbol is used on a type \(\alpha\) ;\\
\item
  mathematical structures, e.g.~the type ring \(\alpha\) that equips a
  type \(\alpha\) with a ring structure;\\
\item
  coercions, via the type class has coe \(\alpha\) \(\beta\) , a wrapper
  around the type \(\alpha \to \beta\) that allows Lean to accept an
  element of \(\alpha\) where a \(\beta\) is expected; and\\
\item
  decidability of propositions, via the decidable p type class, which
  enables case analysis in constructive contexts and if statements in
  programming.
\end{itemize}

This usage is significantly expanded into all parts of mathlib.

Mathematicians and computer scientists often speak of an algebraic
hierarchy, but insofar as that brings to mind the groups, rings, and
fields of an introductory course in abstract algebra, the phrase has the
wrong connotations. Lean does have a hierarchy of algebraic structures,
but this only scratches the surface. For example, in mathlib, the real
numbers are an instance of a normed space, a metric space, a uniform
space, and a normal topological space. Morphisms between structures have
algebraic properties that also need to be managed. The library
instantiates the category of groups and morphisms between them, and the
functor mapping a measure space X to the space of probability measures
on X is an instance of a monad, namely, the Giry monad {[}25{]}. To
convey the richness of such a network of definitions, we will refer to
it rather as a structure hierarchy.

\section{4.1 Organizing the Structure
Hierarchy}\label{organizing-the-structure-hierarchy}

A fragment of the structure hierarchy in Figure 1 shows the classes that
can be derived from a normed\_field instance, which appears near the top
of the hierarchy. The bottom of the graph is populated by notation
classes such as has\_add, has\_le, has\_one, etc. The fragment omits
structures such as partial\_order and lattice, as well as ordered
structures like ordered\_group and ordered\_ring.

Most of the classes in the hierarchy are unary: they have the form
\(C \alpha\) where \(\alpha\) is the carrier type. All of the classes in
the fragment above are unary, but there are some binary classes, such as
module \(\mathsf{R} \mathsf{M}\) that provides an \(R\) -module
structure on \(M\) . The fragment constitutes less than a third of the
current structure hierarchy, which contains 201 unary classes and 266
instances of the form \(C \alpha \to D \alpha\) , which only change the
class associated to a type.

Except for notation classes, all the classes in mathlib have some
independent interest. Uniform spaces are defined because they unify
theorems of topological groups and metric spaces; extended metric spaces
exist because they appear in constructions such as the unbounded
\(\ell^{p}\) spaces. This does not hold for some core library classes,
e.g.~zero\_ne\_one\_class, which exist only to be mixins for other
classes.

Diamonds are common---the structure hierarchy is far from being a tree.
Some of these diamonds arise naturally from the mathematics, some arise
when a class is defined as a least upper bound of two other classes
(such as comm\_monoid), and others arise as a result of encoding a
canonical projection as a parent class.

In some cases, we have a situation in which one structure canonically
implies another: for example, a metric space is not generally considered
to contain a topology as a subcomponent, but the metric uniquely defines
a particular topology of interest. We opt to have metric\_space extend
topological\_space, with an additional constraint asserting that the
topology agrees with the induced topology of the metric. This is done
because Lean depends on definitional equality of projections and this
approach makes more squares commute definitionally. Using this approach,
the induced topology on a product metric becomes definitionally equal to
the product topology on the induced metrics.

\textit{[Image omitted: images/d9c268ff8d7058f1a8927b329ff29a4cdabefb9f8747780cf98196d6da1f30c0.jpg]}

flowchart

Complex community partitioning diagram showing relationships between
normed\_field, normed\_ring, discrete\_field, and various spatial and
community types like uniform\_add\_group, topologicalanding, and
add\_group.

Figure 1. The structure hierarchy underlying normed\_field. An edge from
S to T indicates that an instance of T can be derived from an instance
of S. Some arrows to the leaves are omitted for readability (e.g.~from
zero\_ne\_one\_class to has\_one).

\section{4.1.1 Bundled Type Classes}\label{bundled-type-classes}

When creating a type class, one important design decision is which parts
of the definition to put as parameters to the type class and which parts
to store within the element itself, accessible via a projection. We
refer to a type class as unbundled if it has many parameters and bundled
if the parameters are moved to projections. For example, given the
definitions

\begin{enumerate}
\def\labelenumi{\arabic{enumi}.}
\tightlist
\item
  Group = \{(X, o) \textbar{} (X, o) is a group\}\\
\item
  group \(X = \{\circ \mid (X,\circ)\) is a group\}\\
\item
  is\_group X \(\circ\) \(\leftrightarrow\) (X, \(\circ\) ) is a group
\end{enumerate}

we would call Group a bundled definition, is\_group an unbundled
definition, and group a semi-bundled definition. All of these say
essentially the same thing from a mathematical point of view, but the
choice of which to use has a significant impact on formalization.

We primarily use semi-bundled definitions for type classes in mathlib,
with all operations being bundled except for the carrier types. For use
with type classes, fully bundling (as with Group) is not an option: a
type class problem should have (essentially) at most one solution, and
only the parameters affect the type class search. When the type is
exposed but not the operation, as with group X, we can register a
canonical group associated to the type X if there is one. For instance,
add\_group Z will find the canonical additive group \((\mathbb{Z}, + )\)
of addition on integers. By contrast, fully bundled definitions work
best when using canonical structures, as in the Mathematical Components
library {[}46{]}.

The drawback of further unbundling, as in is\_group X o, is that it
makes matching problems difficult. For example, a theorem saying that f
xy = f y x under the condition is\_comm\_group X f, with x, y, f all
variables, leads to a higher-order matching problem: f could be
substituted for a lambda term. This kind of theorem will always be
applicable, since essentially any term can be written in the form f xy
for some choice of f.~Especially when the simplification tactic simp is
invoked, this can cause a significant performance problem, with many
false positives and failed type class searches. With the partially
unbundled approach this statement instead has the form add X m xy = add
X m y x where m : comm\_group X. We need only look for terms which have
a literal appearance of the constant add in them rather than considering
all terms and excluding them after the fact by a failed type class
search.

\section{4.1.2 Bundled Morphisms}\label{bundled-morphisms}

A similar distinction arises when talking about morphisms between
structures. In most cases, a morphism will be a function with an
additional property, for example a continuous function or a group
homomorphism. In these cases we have

two options: to bundle the function with the property, producing a type
\(\operatorname{Hom}(A, B)\) of structure-respecting functions from
\(A\) to \(B\) , or to have a function \(f: A \to B\) and a type class
is\_hom \(f\) asserting that \(f\) respects the structure of \(A\) and
\(B\) .

Both options are workable, and mathlib contains traces of both
approaches. However, we have found the bundled approach to behave better
for a number of reasons:

\begin{itemize}
\tightlist
\item
  It is important for compositions of homs to be a hom, but problems of
  the form is\_hom ( \(f \circ g\) ) can be difficult for type class
  inference. With bundled homs, one instead defines a custom composition
  operator
  \(\operatorname{Hom}(A, B) \times \operatorname{Hom}(B, C) \to \operatorname{Hom}(A, C)\)
  that is used explicitly, obviating the need to search for is\_hom (
  \(f \circ g\) ).\\
\item
  The type \(\operatorname{Hom}(A, B)\) itself often has interesting
  structure. For example, the type
  \(\mathsf{M} \to_{l}[\mathsf{R}] \mathsf{N}\) of \(R\) -linear maps
  from \(M\) to \(N\) is an \(R\) -module where the zero element is the
  constant zero map, and addition and scalar multiplication work
  pointwise (Section 5.1). Bundling the type makes it easier to reason
  about this structure.
\end{itemize}

\section{4.2 The decidable Type Class}\label{the-decidable-type-class}

The class decidable p is defined essentially as the disjunctive type
\(p \oplus \neg p\) : it is a constructive, Type-valued witness to the
law of excluded middle at proposition p.~Although mathlib primarily
deals in classical mathematics, it is useful to track decidability
because it can be used in algorithms and for small scale computation in
the executable fragment.

The notation if p then t else e is syntactic sugar for ite p t e, where:

\begin{Shaded}
\begin{Highlighting}[]
\NormalTok{ite : ∀ (p : Prop) [d : decidable p],}
\NormalTok{∀ \{α : Sort u\}, α → α → α }
\end{Highlighting}
\end{Shaded}

This function is defined by case analysis on the witness d.~In Lean
syntax, parentheses () denote explicit arguments, curly brackets \{\}
implicit arguments, and square brackets {[}{]} arguments to be inferred
by type class resolution. Notice that p is an explicit argument, but p,
being a proposition, is erased during execution, with the real condition
in d.~Morally, we can think of d as a boolean value, but it carries
evidence with it connecting it to the truth or falsity of p.

The Mathematical Components library {[}46{]} relies heavily on the
technique of small-scale reflection (from which the tactic language
SSReflect gets its name), where many predicates and propositions are
defined to live in type bool. The computational behavior of the logic
can be used to discharge certain proof obligations by reflection. This
means that most properties have two versions, one that is a Prop and one
which is a bool, with an additional predicate reflect b p asserting that
the boolean value b is true iff p holds.

A proof of decidable p is equivalent to \(\Sigma\) b, reflect b p in
this sense, but because it is a type class the management of this side
condition is almost entirely automatic. This means we can treat the
proposition p as primary, as in the definition of ite, and need not
define two versions of every proposition. We can still prove true
decidable propositions by computational reflection via dec\_trivial:

\begin{Shaded}
\begin{Highlighting}[]
\NormalTok{def as\_true (c : Prop) [decidable c] : Prop :=}
\NormalTok{if c then true else false }
\end{Highlighting}
\end{Shaded}

\begin{Shaded}
\begin{Highlighting}[]
\KeywordTok{def}\NormalTok{ of\_as\_true \{c : Prop\} [decidable c] :}
\NormalTok{    as\_true c → c }\OperatorTok{:=}\NormalTok{ ... }
\end{Highlighting}
\end{Shaded}

\begin{Shaded}
\begin{Highlighting}[]
\NormalTok{notation \textasciigrave{}dec\_trivial\textasciigrave{} := of\_as\_true trivial }
\end{Highlighting}
\end{Shaded}

The considerations in this section apply equally to proving,
programming, and metaprogramming in Lean. All propositions can be shown
to be (noncomputably) decidable with the choice axiom; this instance is
used locally with low priority in mathlib to use declarations with
decidability hypotheses in classical mathematics.

\section{4.3 Problems with Type
Classes}\label{problems-with-type-classes}

Given all the aforementioned uses of type classes, and the fact that
most of instances are available in most of the higher-level files, the
growth of mathlib has tested the scalability of the type class inference
system. A file that imports all of mathlib has access to 4256 instances
among 386 classes. While it has held up remarkably well, aspects of the
inference algorithm have lead to limitations in some areas and
performance walls in others.

First, Lean performs a backtracking search on every type class problem.
Thus, an instance such as
\(\mathsf{C} \alpha \leftarrow \mathsf{C} \alpha\) (that is, ``to obtain
\(\mathsf{C} \alpha\) it suffices to prove \(\mathsf{C} \alpha\) ) will
cause the inference routine to enter a loop, blocking resolution even if
a successful path exists elsewhere. As a result, mathlib requires the
complete instance graph to be acyclic. This is a global problem but
generally single instances can be identified as the culprit when an
instance loop is discovered.

Second, and relatedly, because Lean cannot tell when it is retreading
the same paths, even if the graph of instances has no cycles, it will
traverse all paths through the graph, which is exponential time in the
worst case. (On successes it is often significantly less if it gets
lucky, but failures will have to search the whole graph, and
pathological examples can be constructed in which the search is
successful but still exponential time.) Instance searches can take
hundreds of thousands of steps through our thousands of instances.

Third, instance searches are performed exclusively backward, reasoning
from the goal back through instances. This means that if for example G :
group \(\alpha\) is in the context, a search for has\_mul \(\alpha\) may
end up exploring upward through the entire structure hierarchy, looking
for monoid \(\alpha\) , then ring \(\alpha\) , then field \(\alpha\) ,
then normed\_field \(\alpha\) , prompting a search for has\_norm
\(\alpha\) and so on, before eventually failing and attempting group
\(\alpha\) instead. This search will get larger as our structure
hierarchy grows.

While the first two problems seem dire, a combination of good caching
and a very efficient C++ implementation have helped to keep costs
manageable even at mathlib's scale.

Moreover, efficient graph traversal algorithms are known, and Lean 4 is
expected to make improvements in this area.

The third problem could potentially be fixed by using a combination of
forward and backward reasoning. If we reason forward from group
\(\alpha\) to monoid \(\alpha\) , semigroup \(\alpha\) , has\_mul
\(\alpha\) , then work backward from the goal has\_mul \(\alpha\) , we
find the solution quickly without exploring the entire structure
hierarchy. Working forward from instances such as ring \(\mathbb{Z}\)
entails generating additional instances for monoid \(\mathbb{Z}\) ,
add\_group \(\mathbb{Z}\) , and so on, which is currently not required
but is often done to speed up instance searches.

\section{5 Linear Algebra}\label{linear-algebra}

The development of linear algebra in mathlib showcases many of the
design principles explained so far. To make earlier discussions more
concrete, we describe in detail part of the linear algebra development.
This description does not encompass all of the linear algebra in
mathlib: the library also develops the theories of dual vector spaces,
bilinear and sesquilinear forms, direct sums, and tensor products. This
theory contains structures and theorems that are mathematically
nontrivial and are used to prove noteworthy results (Section 7.2). Some
of the structures are close to the top of the type class hierarchy, and
they illustrate the use of bundled and semi-bundled type classes and
quotient types.

Our formalization is certainly not unique: we have been heavily inspired
by existing linear algebra developments, particularly in Isabelle/HOL
{[}3, 21, 43{]} and Coq {[}26{]}.

\section{5.1 Modules}\label{modules}

The fundamental structures of linear algebra are modules. Given a ring
R, an R-module consists of an abelian group M and a distributive,
associative scalar multiplication operator \(\cdot: R \times M \to M\) .
In mathlib, the type class module R M defines an R-module structure on
the group M.

\begin{Shaded}
\begin{Highlighting}[]
\NormalTok{class mul\_action (α : Type u) (β : Type v)}
\NormalTok{    [monoid α] extends has\_scalar α β :=}
\NormalTok{(one\_smul : ∀ b : β, (1 : α) · b = b)}
\NormalTok{(mul\_smul : ∀ (x y : α) (b : β),}
\NormalTok{    (x * y) · b = x · y · b)}

\NormalTok{class distrib\_mul\_action (α : Type u)}
\NormalTok{    (β : Type v) [monoid α] [add\_monoid β]}
\NormalTok{    extends mul\_action α β :=}
\NormalTok{(smul\_add : ∀ (r : α) (x y : β),}
\NormalTok{    r · (x + y) = r · x + r · y)}
\NormalTok{(smul\_zero : ∀ (r : α), r · (0 : β) = 0)}

\NormalTok{class semimodule (R : Type u) (M : Type v)}
\NormalTok{    [semiring R] [add\_comm\_monoid M]}
\NormalTok{    extends distrib\_mul\_action R M :=}
\NormalTok{(add\_smul : ∀ (r s : R) (x : M),}
\NormalTok{    (r + s) · x = r · x + s · x)}
\NormalTok{(zero\_smul : ∀x : M, (0 : R) · x = 0) }
\end{Highlighting}
\end{Shaded}

\begin{Shaded}
\begin{Highlighting}[]
\KeywordTok{class} \KeywordTok{module}\NormalTok{ (R : }\DataTypeTok{Type}\NormalTok{ u) (M : }\DataTypeTok{Type}\NormalTok{ v)}
\NormalTok{    [ring R] [add\_comm\_group M]}
\NormalTok{    extends semimodule R M }
\end{Highlighting}
\end{Shaded}

We provide an alternative constructor for module R M that does not
require the fields smul\_zero and zero\_smul, since they follow from
other properties.

When R is replaced with a field K, this structure is called a K-vector
space. All results about modules also apply to vector spaces. In
mathlib, we favor working with discrete fields, which fix 1/0 = 0.

\begin{Shaded}
\begin{Highlighting}[]
\NormalTok{class vector\_space (K : Type u) (V : Type v)}
\NormalTok{[discrete\_field K] [add\_comm\_group V]}
\NormalTok{extends module K V }
\end{Highlighting}
\end{Shaded}

We consider R-modules for the rest of this section. For the sake of
display, we will omit the following parameters to all declarations:

\begin{Shaded}
\begin{Highlighting}[]
\NormalTok{(R : }\DataTypeTok{Type}\NormalTok{ u) (M : }\DataTypeTok{Type}\NormalTok{ v) [ring R]}
\NormalTok{[add\_comm\_group M] [}\KeywordTok{module}\NormalTok{ R M] }
\end{Highlighting}
\end{Shaded}

We will also sometimes use ``\ldots'' to elide proof terms.

When \((M,\cdot)\) is an R-module, a subgroup of M closed under
\(\cdot\) is called a submodule. These are also defined as bundled
structures in mathlib. Ordered by set inclusion, the submodules of a
module form a complete lattice.

\begin{Shaded}
\begin{Highlighting}[]
\NormalTok{structure submodule :=}
\NormalTok{(carrier : set M)}
\NormalTok{(zero : (0:M) ∈ carrier)}
\NormalTok{(add : ∀ \{x y\}, x ∈ carrier →}
\NormalTok{    y ∈ carrier → x + y ∈ carrier)}
\NormalTok{(smul : ∀ (c:R) \{x\},}
\NormalTok{    x ∈ carrier → c · x ∈ carrier) }
\end{Highlighting}
\end{Shaded}

The lattice structure on submodules allows us to use the notation
\(\perp\) for the trivial submodule 0 and \(\top\) for the universal
submodule M. We also use the lattice structure to define the span of a
set, the space of all linear combinations of its elements:

\begin{Shaded}
\begin{Highlighting}[]
\NormalTok{def span (s : set M) : submodule R M := Inf \{p | s ⊆ p\} }
\end{Highlighting}
\end{Shaded}

The quotient M/N of an R-module M by a submodule \(N \subseteq M\)
equates \(m_{1}, m_{2} \in M\) if \(m_{1} - m_{2} \in N\) . The quotient
is itself an R-module, and is defined in mathlib as a quotient type.

\begin{Shaded}
\begin{Highlighting}[]
\NormalTok{def quotient\_rel (N : submodule R M) :}
\NormalTok{    setoid M := ⟨λ x y, x {-} y ∈ N, ...⟩}
\NormalTok{def submodule.quotient (N : submodule R M) :}
\NormalTok{    Type v := quotient (quotient\_rel N)}
\NormalTok{instance (N : submodule R M) :}
\NormalTok{    module R (quotient N) := ... }
\end{Highlighting}
\end{Shaded}

A function between two R-modules that preserves addition and scalar
multiplication is called a linear map. If it is invertible, it is called
a linear equivalence. We work with linear maps as bundled structures in
mathlib, and coercions

typically allow us to treat these structures as functions. We use the
notation \(M \rightarrow_{l}[R]\) N to stand for linear\_map R M N and
\(M \simeq_{l}[R]\) N for linear\_equiv R M N.

\begin{Shaded}
\begin{Highlighting}[]
\NormalTok{structure linear\_map (N : Type w)}
\NormalTok{[add\_comm\_group N] [module R N] :=}
\NormalTok{(to\_fun : M → N)}
\NormalTok{(add : ∀ x y,}
\NormalTok{    to\_fun (x + y) = to\_fun x + to\_fun y)}
\NormalTok{(smul : ∀ (c : R) x,}
\NormalTok{    to\_fun (c · x) = c · to\_fun x) }
\end{Highlighting}
\end{Shaded}

The type M →l{[}R{]} N is itself an R-module.

If \(f: M \to N\) is linear, the image of a submodule \(M' \subseteq M\)
under \(f\) is a submodule of \(N\) , and the preimage of a submodule
\(N' \subseteq M\) under \(f\) is a submodule of \(M\) . These are
defined respectively as map and comap. In particular, the kernel of
\(f\) -defined to be the set \(\{x \in M \mid f(x) = 0\}\) -is a
submodule, as is the range.

def ker (f : M → \(_{l}\) {[}R{]} N) : submodule R M := comap f ⊥

\begin{Shaded}
\begin{Highlighting}[]
\NormalTok{def range (f : M →l[R] N) : submodule R N := map f T }
\end{Highlighting}
\end{Shaded}

These definitions and their surrounding proofs are enough to prove the
first and second isomorphism laws for modules:

def quot\_ker\_equiv\_range : f.ker.quotient \(\simeq_{l}[R]\) f.range
:= \ldots{}

\begin{Shaded}
\begin{Highlighting}[]
\NormalTok{def sup\_quotient\_equiv\_quotient\_inf}
\NormalTok{(p p\textquotesingle{} : submodule R M) :}
\NormalTok{(comap p.subtype (p ∩ p\textquotesingle{})).quotient ≈l[R]}
\NormalTok{(comap (p ∩ p\textquotesingle{}).subtype p\textquotesingle{}).quotient := ... }
\end{Highlighting}
\end{Shaded}

\section{5.2 Bases}\label{bases}

A set of elements of an \(R\) -module \(\{v_i\}_{i \in I}\) is said to
be linearly dependent if there is a finite subset \(I' \subseteq I\) and
a family of scalars \(\{c_i\}_{i \in I'}\) , not all zero, such that
\(\sum_{i \in I'} c_i \cdot v_i = 0\) . We are more often interested in
the negation of this notion. In mathlib, we define the linear
independence of a family of vectors indexed by a base type:

def linear\_independent

\begin{Shaded}
\begin{Highlighting}[]
\NormalTok{\{ι : Type u\} (R : Type v) \{M : Type w\}}
\NormalTok{[ring R] [add\_comm\_group M] [module R M]}
\NormalTok{(v : ι → M) : Prop :=}
\NormalTok{finsupp.total ι M R v).ker = ⊥ }
\end{Highlighting}
\end{Shaded}

The function finsupp.total \(\iota\) M R v:
\((\iota \to_{0} R) \to_{l}[R]\) M sends a finitely supported function
c: \(\iota \to R\) to the sum over \(\iota\) of c i \(\cdot\) v i.

Through the rest of this subsection we fix parameters as given to
linear\_independent. A linearly independent family of vectors is a basis
for a module if it spans the entire module.

def is\_basis (v : \(\iota \rightarrow M\) ) : Prop :=
linear\_independent R v \(\wedge\) span R (range v) = T

Given such a v and a proof hv : is\_basis R v, we define the linear map
that decomposes a vector into a weighted sum of vectors in v:

\begin{Shaded}
\begin{Highlighting}[]
\NormalTok{def is\_basis.repr : M →l (l →0 R) :=}
\NormalTok{hv.1.repr.comp}
\NormalTok{(linear\_map.id.cod\_restrict \_ hv.mem\_span) }
\end{Highlighting}
\end{Shaded}

A classic result in linear algebra shows that every vector space has a
basis. (This result uses Zorn's lemma and does not hold constructively.)

\begin{Shaded}
\begin{Highlighting}[]
\NormalTok{lemma exists\_is\_basis [discrete\_field K]}
\NormalTok{[add\_comm\_group V] [vector\_space K V] :}
\NormalTok{∃b : set V, is\_basis K (λ i : b, i.val) := ... }
\end{Highlighting}
\end{Shaded}

\section{5.3 Dimension}\label{dimension}

Any two bases for a vector space V have equal cardinality. This
cardinality defines the dimension of V. Vector spaces are not
necessarily finite dimensional, so we define the dimension function to
take values in the type cardinal. The theory of cardinal numbers in
mathlib is located in the set\_theory subfolder. To avoid inconsistency,
the type cardinal must be parametrized by a universe level.

\begin{Shaded}
\begin{Highlighting}[]
\NormalTok{def vector\_space.dim (K : Type u) (V : Type v)}
\NormalTok{    [discrete\_field K] [add\_comm\_group V]}
\NormalTok{    [vector\_space K V] : cardinal.\{v\} :=}
\NormalTok{cardinal.min (nonempty\_subtype.2}
\NormalTok{(@exists\_is\_basis K V \_ \_ \_))}
\NormalTok{(λ b, cardinal.mk b.1) }
\end{Highlighting}
\end{Shaded}

Because the relevant cardinals may live in different universes, it takes
some care to state the dimension theorem.

\begin{Shaded}
\begin{Highlighting}[]
\NormalTok{theorem mk\_eq\_mk\_of\_basis}
\NormalTok{\{v : l → V\} \{v\textquotesingle{} : l\textquotesingle{} → V\}}
\NormalTok{(hv : is\_basis K v) (hv\textquotesingle{} : is\_basis K v\textquotesingle{}) :}
\NormalTok{cardinal\_lift.\{w w\textquotesingle{}\} (cardinal.mk l) =}
\NormalTok{cardinal\_lift.\{w\textquotesingle{} w\} (cardinal.mk l\textquotesingle{}) := ... }
\end{Highlighting}
\end{Shaded}

Many dimension computations are proved in mathlib, including results
about the dimensions of quotients and the rank-nullity theorem.

\begin{Shaded}
\begin{Highlighting}[]
\NormalTok{theorem dim\_quotient (V\textquotesingle{} : submodule K V) :}
\NormalTok{dim K p.quotient + dim K V\textquotesingle{} = dim K V := ... }
\end{Highlighting}
\end{Shaded}

\begin{Shaded}
\begin{Highlighting}[]
\NormalTok{theorem dim\_range\_add\_dim\_ker (f : V → l[K] E) :}
\NormalTok{dim K f.}\BuiltInTok{range} \OperatorTok{+}\NormalTok{ dim K f.ker }\OperatorTok{=}\NormalTok{ dim K V }\OperatorTok{:=}\NormalTok{ ... }
\end{Highlighting}
\end{Shaded}

A number of these computations specialize to finite-dimensional vector
spaces. For instance, the functions \(A \rightarrow k\) form a k-vector
space with dimension \(|A|\) when A is finite.

\begin{Shaded}
\begin{Highlighting}[]
\NormalTok{lemma dim\_fun [fintype η]:}
\NormalTok{dim K (η → K) = fintype.card η := ... }
\end{Highlighting}
\end{Shaded}

\begin{Shaded}
\begin{Highlighting}[]
\NormalTok{lemma dim\_fin\_fun (n : N) :}
\NormalTok{dim K (fin n → K) = n := ... }
\end{Highlighting}
\end{Shaded}

The vector space fin \(n \rightarrow K\) is a common way to represent
n-tuples of elements of K.

\section{5.4 Matrices}\label{matrices}

It is convenient to consider matrices as functions
\(m \rightarrow n \rightarrow R\) , where m and n are finite sets of
indices. We define this type in the data directory of mathlib, along
with the expected operations and algebraic instances. For a ring R,
R-valued \(m \times n\) matrices form an R-module.

\begin{Shaded}
\begin{Highlighting}[]
\NormalTok{def matrix (m n : Type u) (R : Type v)}
\NormalTok{    [fintype m] [fintype n] : Type (max u v) := m → n → R}
\NormalTok{instance [ring R] : module R (matrix m n R) := ... }
\end{Highlighting}
\end{Shaded}

There is a direct correspondence between R-valued matrices and linear
maps between R-modules. To begin with, every such \(m \times n\) matrix
gives rise to a linear map from \(R^{n}\) to \(R^{m}\) . In fact, this
evaluation function is itself a linear map.

def eval : (matrix m n R) → \(_{l}\) {[}R{]} ((n → R) → \(_{l}\) {[}R{]}
(m → R)) := linear\_map.mk \(_{2}\) R mul\_vec \ldots{}

Similarly, a linear map exists in the opposite direction.

def to\_matrix : ((n → R) → \(_{l}\) {[}R{]} (m → R)) → \(_{l}\) {[}R{]}
matrix m n R := linear\_map.mk (λ f i j, f (λ n, ite (j = n) 1 0) i)
\ldots{}

These linear maps are inverses of each other, and thus form an
equivalence between the space of linear maps \(R^{n} \rightarrow R^{m}\)
and the space of R-valued \(m \times n\) matrices. It follows quickly
that, given finite bases for two R-modules \(M_{1}\) and \(M_{2}\)
indexed by \(\iota\) and \(\kappa\) , the space of linear maps
\(M_{1} \rightarrow M_{2}\) is linearly equivalent to the space of
R-valued \(\iota \times \kappa\) matrices.

\begin{Shaded}
\begin{Highlighting}[]
\NormalTok{def lin\_equiv\_matrix \{l : Type i\} \{k : Type k\}}
\NormalTok{    [fintype l] [decidable\_eq l]}
\NormalTok{    [fintype k] [decidable\_eq k]}
\NormalTok{    \{v1 : l → M1\} (hv1 : is\_basis R v1)}
\NormalTok{    \{v2 : k → M2\} (hv2 : is\_basis R v2) :}
\NormalTok{    (M1 →l[R] M2) ≈l[R] matrix k l R := ... }
\end{Highlighting}
\end{Shaded}

\section{5.5 Summary}\label{summary}

While mathlib's algebraic content is significant, it is by no means the
majority of the library, and we do not mean to imply this by focusing on
linear algebra as an example. The preceding sections intend to give the
reader an idea of the flavor of mathlib formalizations more generally.
The use of type classes to manage algebraic and order structure is
representative of the whole library.

\section{6 Metaprogramming}\label{metaprogramming}

Due to Lean's powerful metaprogramming framework {[}22{]}, many features
that might otherwise require changes to the prover or extension through
plugins are implemented in mathlib itself. These include general purpose
decision procedures, utilities, debugging tools, and special purpose
automation. All of the following tactics are implemented in Lean as
metaprograms, with the exceptions of the core simplification routine,
E-matching {[}19{]}, and congruence closure {[}54{]}, which are exposed
in the metalanguage as atomic procedures.

\section{6.1 Simplification}\label{simplification}

The tactic simp is the primary tool in the Lean automation arsenal.
Similar to the simplifier in Isabelle, it performs non-definitional
directional rewriting with equality and biconditional lemmas. Lean's
attribute mechanism allows the user to annotate definitions and theorems
with extra information, which can be accessed by metaprograms; the
default set of rewrite rules for simp is extended by adding the simp
attribute to declarations. A variant, dsimp, performs only definitional
reductions.

The simplifier matches rewrite lemmas up to syntactic, not definitional,
equality. For this reason, lemmas in mathlib are typically stated in
simp-normal form. When there are multiple equivalent ways to express a
term, one is preferred, and others are simplified to the preferred form;
subsequent lemmas need only be stated for this single case. An example
of this design pattern can be seen in the development of the p-adic
numbers Q\_{[}p{]}, where the generic norm notation
\textbar\textbar x\textbar\textbar{} from the normed\_space type class
is preferred to the definitionally equal padic\_norm p x. We use simp
lemmas to transform the latter to the former, and develop the library
for the p-adic norm based on the generic notation.

\begin{Shaded}
\begin{Highlighting}[]
\NormalTok{@[simp] lemma is\_norm (q : Q\_[p]) :}
\NormalTok{(padic\_norm\_e q) = ||q|| := rfl}

\NormalTok{@[simp] lemma mul (q r : Q\_[p]) :}
\NormalTok{||q * r|| = ||q|| * ||r|| := ... }
\end{Highlighting}
\end{Shaded}

\section{6.2 Tactics and Automation}\label{tactics-and-automation}

Besides simp, we briefly describe notable examples of automation in
mathlib, which we roughly classify according to whether they are ``big''
(general-purpose, powerful, and meant to discharge certain classes of
goals on their own) or ``small'' (specialized and providing finer
control over the proof state). Alongside tactics, we include user
commands, which manipulate the prover environment outside of any
particular proof state, and attributes, which can trigger metaprograms
when applied to declarations.

Big Automation For general-purpose proof search, the most powerful tools
in mathlib are finish and tidy. The former combines a tableau prover
(complete for propositional logic) with simplification, E-matching, and
congruence closure. The latter implements heuristics and repeatedly
attempts to apply a preselected list of tactics (recursing into
subgoals) until none succeed. In the category theory library, tidy is
used extensively to check functoriality and naturality. The constructors
for many structures call tidy by default, so that proofs of these
``obvious'' properties need not even be mentioned.

The tactics ring and abel normalize expressions in commutative
(semi)rings and abelian groups. They follow the approach described by
Grégoire and Mahboubi {[}30{]}, but produce proof terms tracing each
normalization step instead of verifying by reflection.

Two tactics are used for solving linear inequalities. For linear ordered
commutative semirings, linearith implements an algorithm based on
Fourier-Motzkin variable elimination {[}61{]}; it is complete for dense
linear orders. On \(\mathbb{N}\) and \(\mathbb{Z}\) , omega partially
implements the omega decision procedure for Presburger arithmetic {[}6,
53{]}.

Small Automation While Lean's metaprogramming engine is powerful enough
for large tactics, it shines in the context of special-purpose tools,
which are often simple to create and deploy. We highlight here a few of
the dozens of such tactics and commands implemented in mathlib.

The norm\_cast tactic manipulates type coercions {[}45{]}. Because Lean
has no subtyping, a coercion (written with a prefix \(\uparrow\) ) is
required, for instance, when using a natural number in place of an
integer. The presence of these coercions in terms can hinder
simplification and unification: it can be tedious to reduce the integer
inequality \(\uparrow m + \uparrow n > 5\) to the natural number
inequality \(m + n > 5\) . Using attributes to track lemmas that pass
coercions through operations and relations, norm\_cast performs such
simplification automatically.

Arithmetic expressions involving only literals are evaluated efficiently
by norm\_num. Over the natural numbers, goals such as \(1 + 2 < 4\) can
be proved by kernel computation. However, this is inefficient for large
expressions, and impossible on noncomputable types such as R and on
arbitrary linear ordered semirings. As an alternative, norm\_num uses
the syntactic binary representation of numerals and lemmas about the
relevant algebraic structures to simplify the arithmetic expressions as
far as possible.

For a type class C, the tactic pi\_instance generalizes instances of the
form C \(\alpha\) to C ( \(\beta \rightarrow \alpha\) ). With this, we
obtain much of the algebraic structure of function spaces for free.

The reassoc attribute improves the ability of simp to reason about
compositions of arrows in a category modulo associativity. By default,
simp normalizes such expressions by associating composition to the
right. As a consequence, one often encounters series of arrow
compositions that are bracketed as a \(\gg\) (b \(\gg\) c). A lemma foo
that rewrites a \(\gg\) b to d would conflict with this simp normal
form. Applying the attribute reassoc to foo produces a companion lemma
of the shape \(\forall X\) , a \(\gg\) (b \(\gg\) X) = d \(\gg\) X that
can be used by the simplifier.

Maintenance and Fine-Tuning Another group of meta- programs is used for
library development and maintenance.

Some search-based tactics can report traces of a successful search.
Variants of tidy and simp list the tactic script and used lemmas,
respectively, in a format that is faster and more robust than the
original tactic call. In editors that support Lean's hole commands,
calls to these tactics can be literally self-replacing.

The command \#lint is a linting tool that performs various tests on
declarations in a file or environment. Among other things, it checks for
unused arguments, malformed names, and whether a file meets mathlib's
documentation requirements. It is easily extended with custom tests.

Lean's namespaces and sectioning mechanisms allow the use of abbreviated
identifiers, special notation, and instances locally in files, as well
as to automatically insert parameters to declarations within a section.
The \#where command prints information about the currently open
namespaces and parameters; the localized command associates local
notation and instances with a namespace, making it easy to set up a
particular local environment.

\section{6.3 Lessons}\label{lessons}

The various tools listed above are provided as a part of mathlib,
without need for extensions or plugins to the Lean system. (One
exception, namely simp, is built into Lean.)

Collaboration between mathematically-oriented formalizers and more
experienced programmers has driven much of the tactic development in
mathlib. Tools such as norm\_num, ring, and norm\_cast arose after user
requests to automate mathematically trivial proofs. Others, such as
pi\_instance, make clever use of the metaprogramming API to eliminate
boilerplate code. It is hard to separate the metaprogram components of
mathlib from the mathematical formalizations, as many tactics were
inspired by particular patterns of use and rely on nontrivial theories
or the proper use of attributes on library lemmas.

While our approach is not unique, our emphasis is a shift from the
traditional point of view: we consider tactic and tool development as
part of library design rather than system development. This division of
labor has been possible thanks to the flexibility of metaprogramming in
Lean.

\section{7 Community}\label{community}

The mathlib library is not developed to support any single project. The
interests of its contributors range from research mathematics, to STEM
education, to automated proof search, to program verification. These
contributors come primarily from academia; some are renowned researchers
and some are bachelor students. Design decisions and future directions
are discussed openly and publicly, with little central control over the
library's content.

That a cohesive library has been developed by a community with such
diverse motivations and backgrounds is perhaps surprising. Unlike most
proof assistant libraries, which are developed or overseen by dedicated
research groups, mathlib is organized as an open-source community
project. While we have witnessed many of the familiar difficulties

with such projects, this organizational scheme is arguably one of the
reasons the library has been successful.

The communal nature of the project is, perhaps, one reason that Lean and
mathlib have attracted many mathematicians. Mathematical topics are
delicately intertwined. Rather than individually building projects on
top of a generic core library, the involved mathematicians have
integrated their projects with each other and with mathlib itself. To
many, the main research question in mathlib is how to design a library
of formal mathematics that is both broad and deep. This question cannot
be answered without the collaboration of many people; the range of
contributors is a vital aspect of the project.

A side effect of the adoption of mathlib by mathematicians is that Lean
is being used in mathematics education. While it is not uncommon to see
proof assistants in certain computer science courses, we know of few
cases where they have been introduced to undergraduate mathematicians
{[}48{]}. Some of these undergraduates interact on the Lean Zulip chat
room, contribute to mathlib, and participate in the review process.

\section{7.1 GitHub and Zulip}\label{github-and-zulip}

Communication about mathlib occurs mainly over two channels. The project
is hosted on GitHub, \(^{2}\) where contributions in the form of pull
requests can be reviewed and edited. More casual conversations occur in
a Zulip chat room. \(^{3}\)

Eleven community members have been designated as maintainers of the
library. These people are responsible for approving pull requests in
their areas of expertise. Pull request reviews are welcomed from all
members of the community; in addition to the formal process on GitHub,
PRs are often discussed on Zulip. Community members are given write
access to non-master branches of the repository, and many PRs are
developed as group efforts on these branches.

The Zulip chat room is home to a broad range of discussions. In
particular, a channel for new members welcomes elementary questions
about Lean and formalization. Another channel focuses on formalizing
advanced mathematics. A number of people have cited the accessibility of
this chat room as a reason for deciding to use Lean.

\section{7.2 Projects Based on mathlib}\label{projects-based-on-mathlib}

As well as being a cohesive collection of mathematics on its own,
mathlib should also serve as a library on which to build more
specialized projects. Users have undertaken many such projects, and the
variety of topics reflects the diverse interests of the mathlib
community. These projects range from published research papers to
collaborative weekend efforts coordinated on Zulip. We list here a
non-exhaustive selection, to give an idea of what mathlib can support.

The cap set problem In 2016, Ellenberg and Gijswijt discovered a
solution to the cap set problem, a longstanding open question in
combinatorics. Their celebrated proof {[}23{]}, published in the Annals
of Mathematics in 2017, was noted for its use of elementary methods from
linear algebra. Dahmen, Hölzl, and Lewis {[}18{]} formalized this result
in 2019. The formalization is based on mathlib and resulted in
significant contributions to the linear algebra theory (Section 5), as
well as those related to finite combinatorics. It is a rare example of
contemporary research mathematics being formalized in a proof assistant,
made possible by collaboration between a mathematician and experts in
formalization.

The continuum hypothesis Han and van Doorn {[}33{]} verified the
unprovability in ZFC of the continuum hypothesis. They have since shown
the unprovability of \(\neg\) CH {[}34{]}, completing the notoriously
intricate full independence proof. The formalization builds on the
mathlib embedding of ZFC and develops the syntactic theory of first
order logic. The independence of CH was one of the few remaining results
on Wiedijk's list of formalization targets \(^{4}\) that had yet to be
formalized in any proof assistant.

Perfectoid spaces Scholze was awarded a Fields Medal in 2018, in part
for introducing the definition of a perfectoid space. To test the
ability of a proof assistant to understand extremely complicated
mathematical structures, Buzzard, Commelin, and Massot defined
perfectoid spaces in Lean {[}15{]}. Defining this structure relies on a
large amount of algebraic and topological theory, and the months-long
effort to formalize it resulted in thousands of lines of Lean code.

The sensitivity conjecture Huang's sensational proof of the boolean
sensitivity conjecture {[}37{]} in 2019 was widely discussed, including
on the Lean Zulip chat. Following Knuth's simplified writeup, \(^{5}\) a
number of community members formalized the proof in a matter of days.
The formalization reuses linear algebra machinery from the cap set
project and amounts to fewer than 450 lines of heavily commented code.

Cubing a cube After a challenge was issued at the 2019 Big Proofs
meeting in Edinburgh, van Doorn formalized J. E. Littlewood's
``elegant'' proof that any dissection of a cube into smaller cubes must
contain at least two cubes of equal width. This had been considered an
exemplar of a proof with a simple intuitive argument that is hard to
make formal. It was another remaining item on Wiedijk's list of targets.

Decision procedures for modal logics Wu and Goré {[}62{]} verified
decision procedures for the modal logics K, KT, and S4. The decision
procedures are formalized as programs in Lean with total correctness
proved. The formalization makes extensive use of sublist permutation
which is fully supported by mathlib. The need for proving termination
which involves

reasoning about modal degrees has in turn resulted in further
development of list theory of mathlib.

\section{8 Comparison with Other
Libraries}\label{comparison-with-other-libraries}

In this section, we compare and contrast mathlib with other substantial
formal libraries for mathematics, including libraries for Mizar {[}9,
28{]}, HOL Light {[}35{]}, Isabelle/HOL {[}49{]}, Coq/SSReflect {[}10,
46{]}, and Metamath {[}47{]}. Our goal here is not to provide detailed
comparisons of the various design choices, but, rather to sketch the
design space in broad strokes, situate mathlib within it, and explain
some of the decisions we have made. Our choice of comparisons is not
meant to be exhaustive: there are also substantial mathematical
libraries in HOL4 {[}55{]}, ACL2 {[}42{]}, PVS {[}50{]}, and NuPRL
{[}17{]}, as well as notable libraries built on standard Coq, such as
the Coquelicot analysis library {[}14{]}.

\section{8.1 The Design Space}\label{the-design-space}

We will focus on three specific aspects of the design space, namely, the
choice of axiomatic foundation, the mechanisms used to represent
structures and the relationships between them, and the use of
automation.

Most contemporary interactive proof systems are based on either set
theory, simple type theory, or dependent type theory. The decision to
use an untyped framework like set theory or a typed alternative is
fundamental. Types play two important roles. First, input in a typed
system is often less verbose because users can elide details that can be
inferred from the types of the objects involved. For example, users can
write x + y and allow the system to infer the meaning of the plus sign
from the types of its arguments. Second, a type system can more easily
catch and report errors, such as sending the wrong number or kinds of
arguments to a function, or sending them in the wrong order. Another
important choice is whether or not the axiomatic framework is
constructive, and, in particular, whether it specifies a computational
behavior for objects defined in the framework.

\section{8.2 Libraries Based on Set
Theory}\label{libraries-based-on-set-theory}

As an axiomatic foundation, set theory has two important advantages.
First, it is accepted by many mathematicians as being the official
foundation for mathematics, or, at least, an uncontroversial one. And,
second, it is fairly easy to implement and check proofs in this
framework.

The Metamath system {[}47{]} is a generic system for representing formal
axiomatic frameworks, but its largest library is built on set theory.
That library currently has about 23,000 theorems and 600k lines of code,
covering algebra, analysis, topology, number theory, and other areas.
Because the foundation is so simple, the source files need to provide
explicit proofs that formulas are well formed, and there are front ends
that insert that automatically. But the system uses very little
automation beyond that. There are a number of reference checkers on
offer that can check the entire library in a few seconds. Sets are used
to relativize quantifiers to specific domains. For example, the
following states that every continuous function on a closed interval is
bounded:

\[
\begin{array}{r l}&((A \in \mathbb {R} \land B \in \mathbb {R} \land F \in ((A [, ] B) c n \rightarrow \mathbb {C})) \rightarrow\\&\exists x \in \mathbb {R} \forall y \in (A [, ] B) (a b s ^ {\prime} (F ^ {\prime} y)) \leq x)\end{array}
\]

The need to write proofs that are fully detailed and explicit, however,
places a high burden on the user.

Mizar {[}9, 28{]}, which dates back to the early 1970s, is based on set
theory with Grothendieck-Tarski universes. Its vast library {[}8{]}
currently contains over 3.1 million lines of code and spans many fields
of mathematics. Proofs in the system are designed to mirror mathematical
vernacular. They are written in a declarative way, and a fixed checker
determines whether inferences are valid {[}28{]}. To support a structure
hierarchy, Mizar uses a soft typing system whereby users register
associations that are used by the checker. The library has developed
quite an elaborate hierarchy in this way {[}29{]} (see also {[}28,
39{]}). To our knowledge, there is no formal specification of the
inferences that are accepted by the system, making it hard to implement
an independent reference checker. Nonetheless, Kaliszyk et al.~have had
success reimplementing the system in the Isabelle framework and checking
some of the Mizar files {[}39, 40{]}, and Urban and Sutcliffe have shown
that Mizar's inferences can be cross validated by automated theorem
provers {[}58{]}.

\section{8.3 Libraries Based on Simple Type
Theory}\label{libraries-based-on-simple-type-theory}

A number of contemporary proof systems implement versions of simple type
theory, among them HOL4 {[}55{]}, HOL Light {[}35{]}, and Isabelle/HOL
{[}49{]}. In simple type theory, types and objects are separate things;
one defines types and type constructions, and then one defines objects
of those types. Types cannot depend on objects: for example, one cannot
define a type \(\mathbb{R}^n\) that depends on a natural number
parameter \(n\) .

Simple type theory excels at dealing with concrete structures like the
integers, the reals, and the complex numbers, but when it comes to
algebraic reasoning, the fact that types cannot depend on parameters is
a severe restriction. Structures in mathematics are often parameterized
by other objects: for \(n \geq 1\) , the type of \(M_{n}(R)\) of
\(n \times n\) over a ring form a ring, and for every prime number p,
the integers modulo p, Z/pZ, form a field; \(L^{p}\) spaces, rings of
p-adic integers, and spaces \(C^{k}(X)\) consisting of k-times
continuously differentiable real-valued functions on a subset X of the
reals all depend on parameters. In simple type theory, one has to model
these using predicates on fixed ambient types, erasing many of the
benefits of using type theory in the first place. Another limitation is
that one cannot form spaces or structures whose

elements are structures, such as the space of nonempty compact metric
spaces with the Gromov-Hausdorff distance or the category of groups in
some universe.

HOL Light {[}35{]} is John Harrison's implementation of simple type
theory, inspired by Mike Gordon's original HOL system. The library,
close to 800k lines of code, is especially strong in multivariate real
analysis and complex analysis. It served as the basis for the Flyspeck
project {[}32{]}. Harrison uses a trick {[}36{]} to formalize theorems
about \(\mathbb{R}^n\) for arbitrary \(n\) , using type variables to
stand proxy for their cardinality. But this trick has limited utility:
in simple type theory, one cannot even state that for every \(n\) there
is a type of cardinality \(n\) , and so there is no way of instantiating
such a generic theorem to particular parameters.

Many aspects of our structure hierarchy, including the algebraic
hierarchy, the axiomatization of topological spaces and normed spaces,
our development of measure theory, and our use of filters in analysis,
were modeled after similar developments in Isabelle {[}49{]}. Isabelle's
extensive library has about 900k lines of code, and its Archive of
Formal Proofs {[}12{]} has an additional 2.3 million lines of code. To
treat common data types as instances of structures, Isabelle adopts a
conservative extension of simple type theory with axiomatic type classes
{[}31, 60{]}. But type classes can only depend on a single type
parameter, and, moreover, the use of type classes suffers from the
limitations imposed by simple types. For example, the extensive
libraries and automation developed for instances of rings like the
integers and the reals cannot be applied to the ring of integers modulo
some number m. For such structures, Isabelle also provides the mechanism
of locales {[}7, 41{]}, but, once again, this means relinquishing some
of the benefits of the use of type theory in the first place. We
therefore consider it a substantial accomplishment that we are beginning
to approximate the depth and range of Isabelle's structure hierarchy in
a dependently typed setting.

Isabelle's supporting automation is especially good. The system provides
internal automation, decision procedures, and a conditional term
rewriter {[}49{]}, but can also call external resolution provers and SMT
solvers and reconstruct their proofs {[}11, 13{]}. Using axiomatic type
classes, internal automation can work generically with types that
instantiate the relevant structures, just as our norm\_num works for
types that support numerals and the relevant operations. Isabelle's
automation is much more powerful than that of Lean and mathlib, but we
again consider it notable that we are making progress towards the use of
such automation with dependent types.

\section{8.4 Libraries Based on Dependent Type
Theory}\label{libraries-based-on-dependent-type-theory}

Turning to dependent type theory, the best point of comparison for
mathlib is the Mathematical Components library {[}46{]}, based on Coq
and the SSReflect proof language {[}27{]}. Coq's version of dependent
type theory is very similar to Lean's. The Mathematical Components
library was the basis for the landmark formalization of the odd order
theorem {[}27{]}, and focuses on related parts of mathematics, including
group theory, linear algebra, and representation theory. The library
itself is about 110k lines of code, not including results specific to
the odd order theorem. Other libraries have been built on it, including
an analysis library {[}2{]} that is still under development.

In contrast to mathlib, the Mathematical Components library is
constructive throughout (although the analysis library just mentioned
uses classical logic). Another distinguishing feature of the
Mathematical Components library is that it uses very little external
automation, and instead relies heavily on the computational
interpretation of the underlying axiomatic framework, so that inferences
can be verified by computational unfolding of expressions.

Like mathlib, Mathematical Components relies on powerful elaboration
mechanisms to support a structure hierarchy, though the specific
mechanisms are different: instead of type classes, Mathematical
Components uses canonical structures {[}24{]}, a particular kind of
unification hint {[}4{]}. The core algebraic hierarchy in the
Mathematical Components library {[}46, Section 7.3{]} includes key
structures such as Z-modules (essentially, abelian groups), rings,
commutative rings, fields, algebraically closed fields, modules, and
algebras. Other parts of the library define vector spaces, and
structures with orders and norms, that are designed to support reasoning
about subfields of the complex numbers. There are additional structures
to support group theory and representation theory, and the analysis
library includes structures such as topological spaces and normed
modules.

The Mathematical Components library is generally more conservative than
we are when it comes to instantiating structures and substructures. For
example, the theory of the natural numbers in Mathematical Components is
largely self contained, whereas our natural numbers instantiate an
ordered semiring, which in turn inherits from the structures of additive
and multiplicative monoids, linear orders, and so on. A calculation
involving the natural numbers in mathlib may, and usually does, involve
generic facts from all levels of this hierarchy.

The Mathematical Components library is also carefully designed to avoid
the need for the kinds of searches described in Section 4. An artifact
of the use of canonical structures is that concrete structures have to
be instantiated to all the classes they inherit from. For example, the
integers are declared as a Z-module, a ring, a commutative ring, an
integral domain, and so on. In mathlib, declaring an instance of a
structure automatically handles not just substructures, but also induced
structure. When we declare the reals to be a metric space, for example,
it thereby inherits the structure of a uniform space hence a topological
space.

\textit{[Image omitted: images/738e0087f296ceb25195c8a9fe61c2aa498194d87a5cb383a8a66f71b3560312.jpg]}

line

{\def\LTcaptype{none} % do not increment counter
\begin{longtable}[]{@{}ll@{}}
\toprule\noalign{}
Date & Lines \\
\midrule\noalign{}
\endhead
\bottomrule\noalign{}
\endlastfoot
2017-07-01 & 0 \\
2017-10-01 & 25000 \\
2018-01-01 & 50000 \\
2018-04-01 & 75000 \\
2018-07-01 & 100000 \\
2018-10-01 & 125000 \\
2019-01-01 & 150000 \\
2019-04-01 & 175000 \\
2019-07-01 & 200000 \\
2019-10-01 & 225000 \\
2020-01-01 & 250000 \\
\end{longtable}
}

Figure 2. Lines of code in the mathlib repository over time, including
white space, comments, and auxiliary files.

\section{9 Conclusion}\label{conclusion}

With respect to the design space described in the previous section, the
characteristic features of mathlib are as follows. First, it is based on
a dependent type theory. We have chosen a typed framework for the
reasons indicated in Section 8.1, and given that we want to carry out
the full range of mathematical constructions, judicious use of dependent
types is unavoidable. Second, it is focused on contemporary mathematics,
which is resolutely classical. Nonetheless, large portions of the
library are explicitly computational, and functions, say, on lists and
integers can be executed. Third, it incorporates useful automation, such
as a good conditional simplifier and domain-specific tactics written in
the system's metaprogramming language. Our automation is still in an
incipient state, but we hope and expect to expand these capabilities.
Finally, it includes a large and interconnected hierarchy of
mathematical structures and instances. Each of the other libraries
discussed in Section 8 shares some of these features, but we know of no
other library that shares them all.

The mathlib library and its surrounding community continue to grow
(Figures 2 and 3). Various individuals and groups are actively working
to expand its mathematical content, automated reasoning tools, and
surrounding infrastructure. There is no target or end goal in mind: as
in traditional mathematics, the development of new theories will follow
the needs and desires of the people involved.

A new version of Lean is currently in development {[}57{]}. It will not
be backward compatible with the current version, and when the system is
ready, we expect to update mathlib. At the current time, it is hard to
predict the scale of this undertaking. Lean 4 promises even more
efficient metaprogramming with access to input syntax trees and a
customizable parser; it might be possible to automate much of the
porting process.

The design of mathlib builds on the lessons of existing libraries, but
the library's special character can be attributed to the range of its
contributors. The combined efforts of mathematicians, computer
scientists, and programmers to design a single unified library of
formalizations and tools has been successful, and we are still learning
from one another, as the community continues to evolve.

\textit{[Image omitted: images/73ea514c68dea66fce6c4400475cb7efb3357e0574065bbfd1282041690a3fa9.jpg]}

line

{\def\LTcaptype{none} % do not increment counter
\begin{longtable}[]{@{}
  >{\raggedright\arraybackslash}p{(\linewidth - 20\tabcolsep) * \real{0.1290}}
  >{\raggedright\arraybackslash}p{(\linewidth - 20\tabcolsep) * \real{0.0860}}
  >{\raggedright\arraybackslash}p{(\linewidth - 20\tabcolsep) * \real{0.0860}}
  >{\raggedright\arraybackslash}p{(\linewidth - 20\tabcolsep) * \real{0.0860}}
  >{\raggedright\arraybackslash}p{(\linewidth - 20\tabcolsep) * \real{0.0860}}
  >{\raggedright\arraybackslash}p{(\linewidth - 20\tabcolsep) * \real{0.0860}}
  >{\raggedright\arraybackslash}p{(\linewidth - 20\tabcolsep) * \real{0.0860}}
  >{\raggedright\arraybackslash}p{(\linewidth - 20\tabcolsep) * \real{0.0860}}
  >{\raggedright\arraybackslash}p{(\linewidth - 20\tabcolsep) * \real{0.0860}}
  >{\raggedright\arraybackslash}p{(\linewidth - 20\tabcolsep) * \real{0.0860}}
  >{\raggedright\arraybackslash}p{(\linewidth - 20\tabcolsep) * \real{0.0968}}@{}}
\toprule\noalign{}
\begin{minipage}[b]{\linewidth}\raggedright
Date
\end{minipage} & \begin{minipage}[b]{\linewidth}\raggedright
Line 1
\end{minipage} & \begin{minipage}[b]{\linewidth}\raggedright
Line 2
\end{minipage} & \begin{minipage}[b]{\linewidth}\raggedright
Line 3
\end{minipage} & \begin{minipage}[b]{\linewidth}\raggedright
Line 4
\end{minipage} & \begin{minipage}[b]{\linewidth}\raggedright
Line 5
\end{minipage} & \begin{minipage}[b]{\linewidth}\raggedright
Line 6
\end{minipage} & \begin{minipage}[b]{\linewidth}\raggedright
Line 7
\end{minipage} & \begin{minipage}[b]{\linewidth}\raggedright
Line 8
\end{minipage} & \begin{minipage}[b]{\linewidth}\raggedright
Line 9
\end{minipage} & \begin{minipage}[b]{\linewidth}\raggedright
Line 10
\end{minipage} \\
\midrule\noalign{}
\endhead
\bottomrule\noalign{}
\endlastfoot
2017-07-01 & 0 & 0 & 0 & 0 & 0 & 0 & 0 & 0 & 0 & 0 \\
2017-10-01 & 50 & 0 & 0 & 0 & 0 & 0 & 0 & 0 & 0 & 0 \\
2018-01-01 & 100 & 0 & 0 & 0 & 0 & 0 & 0 & 0 & 0 & 0 \\
2018-04-01 & 150 & 0 & 0 & 0 & 0 & 0 & 0 & 0 & 0 & 0 \\
2018-07-01 & 200 & 0 & 0 & 0 & 0 & 0 & 0 & 0 & 0 & 0 \\
2018-10-01 & 250 & 0 & 0 & 0 & 0 & 0 & 0 & 0 & 0 & 0 \\
2019-01-01 & 300 & 0 & 0 & 0 & 0 & 0 & 0 & 0 & 0 & 0 \\
2019-04-01 & 350 & 0 & 0 & 0 & 0 & 0 & 0 & 0 & 0 & 0 \\
2019-07-01 & 370 & 0 & 0 & 0 & 0 & 0 & 0 & 0 & 0 & 0 \\
2019-10-01 & 380 & 0 & 0 & 0 & 0 & 0 & 0 & 0 & 0 & 0 \\
2020-01-01 & 380 & 0 & 0 & 0 & 0 & 0 & 0 & 0 & 0 & 0 \\
\end{longtable}
}

Figure 3. Number of commits to the mathlib repository over time, by the
20 most active authors.

\section{Acknowledgments}\label{acknowledgments}

We thank Jasmin Blanchette, Assia Mahboubi, Josef Urban, and the
anonymous referees for their helpful comments and suggestions.

The following people have either authored commits in the mathlib
repository or are attributed authors of mathlib files:

Lucas Allen, Ellen Arlt, Jeremy Avigad, Tim Baanen, Seul Baek, Reid
Barton, Tim Baumann, Alexander Bentkamp, Alex Best, Aaron Bryce, Kevin
Buzzard, Louis Carlin, Mario Carneiro, Bryan Gin-ge Chen, Johan
Commelin, Sander Dahmen, Floris van Doorn, Gabriel Ebner, Ramon
Fernández Mir, Fabian Glöckle, Sébastien Gouëzel, Tobias Grosser, Jesse
Han, Keeley Hoek, Johannes Hölzl, Michael Howes, Simon Hudon,
Christopher Hughes, Michael Jendrusch, Sangwoo Jo, Kevin Kappelmann,
Parikshit Khanna, Yury Kudryashov, Joey van Langen, Kenny Lau, Sean
Leather, Guy Leroy, Robert Y. Lewis, Amelia Livingston, Isabel
Longbottom, Jean Lo, Paul-Nicolas Madelaine, Patrick Massot, Bhavik
Mehta, Rohan Mitta, Stephen Morgan, Scott Morrison, Leonardo de Moura,
Oliver Nash, Wojciech Nawrocki, Morenikeji Neri, Casper Putz, Matt Rice,
Mitchell Rowett, Jan-David Salchow, Blair Shi, Rory Qianyi Shu, Calle
Sönne, Robert Spencer, Neil Strickland, Abhimanyu Pallavi Sudhir, Callum
Sutton, Andreas Swerdlow, Nathaniel Thomas, Alistair Tucker, Sebastian
Ullrich, Koundinya Vajjha, Jens Wagemaker, Minchao Wu, Haitao Zhang,
Zhouhang Zhou, Sebastian Zimmer, Andrew Zipperer.

\section{References}\label{references}

{[}1{]} 1994. The QED Manifesto. In Automated Deduction - CADE-12, 12th
International Conference on Automated Deduction, Nancy, France, June 26
- July 1, 1994, Proceedings. 238--251.
http://link.springer.com/chapter/10.1007/3-540-58156-1\_17\\
{[}2{]} Reynald Affeldt, Cyril Cohen, and Damien Rouhling. 2018.
Formalization Techniques for Asymptotic Reasoning in Classical Analysis.
J.

Formalized Reasoning 11, 1 (2018), 43--76.
https://doi.org/10.6092/issn.1972-5787/8124\\
{[}3{]} Jesús Aransay and Jose Divasón. 2014. Formalization and
Execution of Linear Algebra: From Theorems to Algorithms. In Logic-Based
Program Synthesis and Transformation, Gopal Gupta and Ricardo Peña
(Eds.). Springer International Publishing, Cham, 1--18.
https://doi.org/10.1007/978-3-319-14125-1-1\\
{[}4{]} Andrea Asperti, Wilmer Ricciotti, Claudio Sacerdoti Coen, and
Enrico Tassi. 2009. Hints in Unification. In Theorem Proving in Higher
Order Logics, 22nd International Conference, TPHOLs 2009, Munich,
Germany, August 17-20, 2009. Proceedings. 84--98.
https://doi.org/10.1007/978-3-642-03359-9\_8\\
{[}5{]} Jeremy Avigad, Leonardo de Moura, and Soonho Kong. 2014. Theorem
Proving in Lean. Carnegie Mellon University.\\
{[}6{]} Seulkee Baek. 2019. Reflected Decision Procedures in Lean.
Master's thesis. Carnegie Mellon University.
http://www.andrew.cmu.edu/user/avigad/Students/baek\_ms\_thesis.pdf\\
{[}7{]} Clemens Ballarin. 2003. Locales and Locale Expressions in
Isabelle/Isar. In Types for Proofs and Programs, International Workshop,
TYPES 2003, Torino, Italy, April 30 - May 4, 2003, Revised Selected
Papers. 34--50. https://doi.org/10.1007/978-3-540-24849-1\_3\\
{[}8{]} Grzegorz Bancerek, Czeslaw Bylinski, Adam Grabowski, Artur
Kornilowicz, Roman Matuszewski, Adam Naumowicz, and Karol Pak. 2018. The
Role of the Mizar Mathematical Library for Interactive Proof Development
in Mizar. J. Autom. Reasoning 61, 1-4 (2018), 9--32.
https://doi.org/10.1007/s10817-017-9440-6\\
{[}9{]} Grzegorz Bancerek, Czesław Byliński, Adam Grabowski, Artur
Korniłowicz, Roman Matuszewski, Adam Naumowicz, Karol Pąk, and Josef
Urban. 2015. Mizar: State-of-the-art and Beyond. In Intelligent Computer
Mathematics, Manfred Kerber, Jacques Carette, Cezary Kaliszyk, Florian
Rabe, and Volker Sorge (Eds.). Springer International Publishing, Cham,
261--279. https://doi.org/10.1007/978-3-319-20615-8 17\\
{[}10{]} Yves Bertot and Pierre Castéran. 2004. Interactive Theorem
Proving and Program Development - Coq'Art: The Calculus of Inductive
Constructions. Springer. https://doi.org/10.1007/978-3-662-07964-5\\
{[}11{]} Jasmin Christian Blanchette, Sascha Böhme, and Lawrence C.
Paulson. 2013. Extending Sledgehammer with SMT Solvers. J. Autom.
Reasoning 51, 1 (2013), 109--128.
https://doi.org/10.1007/s10817-013-9278-5\\
{[}12{]} Jasmin Christian Blanchette, Max W. Haslbeck, Daniel Matichuk,
and Tobias Nipkow. 2015. Mining the Archive of Formal Proofs. In
Intelligent Computer Mathematics - International Conference, CICM 2015,
Washington, DC, USA, July 13-17, 2015, Proceedings. 3--17.
https://doi.org/10.1007/978-3-319-20615-8\_1\\
{[}13{]} Jasmin Christian Blanchette, Cezary Kaliszyk, Lawrence C.
Paulson, and Josef Urban. 2016. Hammering towards QED. J. Formalized
Reasoning 9, 1 (2016), 101--148.
https://doi.org/10.6092/issn.1972-5787/4593\\
{[}14{]} Sylvie Boldo, Catherine Lelay, and Guillaume Melquiond. 2015.
Coquelicot: A User-Friendly Library of Real Analysis for Coq.
Mathematics in Computer Science 9, 1 (2015), 41--62.
https://doi.org/10.1007/s11786-014-0181-1\\
{[}15{]} Kevin Buzzard, Johan Commelin, and Patrick Massot. 2020.
Formalising Perfectoid Spaces. In Proceedings of the 9th ACM SIGPLAN
International Conference on Certified Programs and Proofs, CPP 2020, New
Orleans, LA, January 20-21, 2020.\\
{[}16{]} Mario Carneiro. 2019. Formalizing Computability Theory via
Partial Recursive Functions. In 10th International Conference on
Interactive Theorem Proving (ITP 2019) (Leibniz International
Proceedings in Informatics (LIPIcs)), John Harrison, John O'Leary, and
Andrew Tolmach (Eds.), Vol. 141. Schloss Dagstuhl--Leibniz-Zentrum fuer
Informatik, Dagstuhl, Germany, 12:1--12:17.
https://doi.org/10.4230/LIPIcs.ITP.2019.12\\
{[}17{]} Robert L. Constable, Stuart F. Allen, Mark Bromley, Rance
Cleaveland, J. F. Cremer, R. W. Harper, Douglas J. Howe, Todd B.
Knoblock, N. P.

Mendler, Prakash Panangaden, James T. Sasaki, and Scott F. Smith. 1986.
Implementing mathematics with the Nuprl proof development system.
Prentice Hall. http://dl.acm.org/citation.cfm?id=10510\\
{[}18{]} Sander R. Dahmen, Johannes Hölzl, and Robert Y. Lewis. 2019.
Formalizing the Solution to the Cap Set Problem. In 10th International
Conference on Interactive Theorem Proving (ITP 2019) (Leibniz
International Proceedings in Informatics (LIPIcs)), John Harrison, John
O'Leary, and Andrew Tolmach (Eds.), Vol. 141. Schloss
Dagstuhl-Leibniz-Zentrum fuer Informatik, Dagstuhl, Germany,
15:1--15:19. https://doi.org/10.4230/LIPIcs.ITP.2019.15\\
{[}19{]} Leonardo de Moura and Nikolaj Bjørner. 2007. Efficient
E-Matching for SMT Solvers. In Automated Deduction - CADE-21, 21st
International Conference on Automated Deduction, Bremen, Germany, July
17-20, 2007, Proceedings. 183--198.
https://doi.org/10.1007/978-3-540-73595-3\_13\\
{[}20{]} Leonardo de Moura, Soonho Kong, Jeremy Avigad, Floris van
Doorn, and Jakob von Raumer. 2015. The Lean Theorem Prover (system
description). https://leanprover.github.io/papers/system.pdf\\
{[}21{]} Jose Divasón and Jesús Aransay. 2013. Rank-Nullity Theorem in
Linear Algebra. Archive of Formal Proofs (Jan.~2013).
http://isa-afp.org/entries/Rank\_Nullity\_Theorem.html, Formal proof
development.\\
{[}22{]} Gabriel Ebner, Sebastian Ullrich, Jared Roesch, Jeremy Avigad,
and Leonardo de Moura. 2017. A metaprogramming framework for formal
verification. PACMPL 1, ICFP (2017), 34:1--34:29.
https://doi.org/10.1145/3110278\\
{[}23{]} Jordan S. Ellenberg and Dion Gijswijt. 2017. On large subsets
of \(F_{q}^{n}\) with no three-term arithmetic progression. Ann. of
Math. (2) 185, 1 (2017), 339--343.
https://doi.org/10.4007/annals.2017.185.1.8\\
{[}24{]} François Garillot, Georges Gonthier, Assia Mahboubi, and
Laurence Rideau. 2009. Packaging Mathematical Structures. In Theorem
Proving in Higher Order Logics, 22nd International Conference, TPHOLs
2009, Munich, Germany, August 17-20, 2009. Proceedings. 327--342.
https://doi.org/10.1007/978-3-642-03359-9\_23\\
{[}25{]} Michèle Giry. 1982. A categorical approach to probability
theory. In Categorical aspects of topology and analysis (Ottawa, Ont.,
1980). Lecture Notes in Math., Vol. 915. Springer, Berlin-New York,
68--85. https://doi.org/10.1007/BFb0092872\\
{[}26{]} Georges Gonthier. 2011. Point-Free, Set-Free Concrete Linear
Algebra. In Interactive Theorem Proving, Marko van Eekelen, Herman
Geuvers, Julien Schmaltz, and Freek Wiedijk (Eds.). Springer Berlin
Heidelberg, Berlin, Heidelberg, 103--118.
https://doi.org/10.1007/978-3-642-22863-6\_10\\
{[}27{]} Georges Gonthier and Assia Mahboubi. 2010. An introduction to
small scale reflection in Coq. J. Formalized Reasoning 3, 2 (2010),
95--152. https://doi.org/10.6092/issn.1972-5787/1979\\
{[}28{]} Adam Grabowski, Artur Kornilowicz, and Adam Naumowicz. 2010.
Mizar in a Nutshell. Journal of Formalized Reasoning 3, 2 (2010),
153--245. https://doi.org/10.6092/issn.1972-5787/1980\\
{[}29{]} Adam Grabowski, Artur Kornilowicz, and Christoph Schwarzweller.
2016. On algebraic hierarchies in mathematical repository of Mizar. In
Proceedings of the 2016 Federated Conference on Computer Science and
Information Systems, FedCSIS 2016, Gdańsk, Poland, September 11-14,
2016. 363--371. https://doi.org/10.15439/2016F520\\
{[}30{]} Benjamin Grégoire and Assia Mahboubi. 2005. Proving Equalities
in a Commutative Ring Done Right in Coq. In Theorem Proving in Higher
Order Logics, Joe Hurd and Tom Melham (Eds.). Springer Berlin
Heidelberg, Berlin, Heidelberg, 98--113.
https://doi.org/10.1007/11541868\_7\\
{[}31{]} Florian Haftmann and Makarius Wenzel. 2006. Constructive type
classes in Isabelle. In International Workshop on Types for Proofs and
Programs. Springer, 160--174.
https://doi.org/10.1007/978-3-540-74464-1\_11\\
{[}32{]} Thomas C. Hales, Mark Adams, Gertrud Bauer, Dat Tat Dang, John
Harrison, Truong Le Hoang, Cezary Kaliszyk, Victor Magron, Sean
McLaughlin, Thang Tat Nguyen, Truong Quang Nguyen, Tobias Nipkow, Steven
Obua, Joseph Pleso, Jason M. Rute, Alexey Solovyev,

An Hoai Thi Ta, Trung Nam Tran, Diep Thi Trieu, Josef Urban, Ky Khac Vu,
and Roland Zumkeller. 2017. A Formal Proof of the Kepler Conjecture.
Forum of Mathematics, Pi 5 (2017), e2.
https://doi.org/10.1017/fmp.2017.1\\
{[}33{]} Jesse Michael Han and Floris van Doorn. 2019. A Formalization
of Forcing and the Unprovability of the Continuum Hypothesis. In 10th
International Conference on Interactive Theorem Proving (ITP 2019)
(Leibniz International Proceedings in Informatics (LIPIcs)), John
Harrison, John O'Leary, and Andrew Tolmach (Eds.), Vol. 141. Schloss
Dagstuhl--Leibniz-Zentrum fuer Informatik, Dagstuhl, Germany,
19:1--19:19. https://doi.org/10.4230/LIPIcs.ITP.2019.19\\
{[}34{]} Jesse Michael Han and Floris van Doorn. 2020. A Formal Proof of
the Independence of the Continuum Hypothesis. In Proceedings of the 9th
ACM SIGPLAN International Conference on Certified Programs and Proofs,
CPP 2020, New Orleans, LA, January 20-21, 2020.\\
{[}35{]} John Harrison. 2009. HOL Light: An Overview. In Theorem Proving
in Higher Order Logics, Stefan Berghofer, Tobias Nipkow, Christian
Urban, and Makarius Wenzel (Eds.). Springer Berlin Heidelberg, Berlin,
Heidelberg, 60--66. https://doi.org/10.1007/978-3-642-03359-9\_4\\
{[}36{]} John Harrison. 2013. The HOL Light Theory of Euclidean Space.
J. Autom. Reasoning 50, 2 (2013), 173--190.
https://doi.org/10.1007/s10817-012-9250-9\\
{[}37{]} Hao Huang. 2019. Induced subgraphs of hypercubes and a proof of
the Sensitivity Conjecture. arXiv preprint arXiv:1907.00847 (2019).\\
{[}38{]} Fabian Immler and Bohua Zhan. 2019. Smooth manifolds and types
to sets for linear algebra in Isabelle/HOL. In Proceedings of the 8th
ACM SIGPLAN International Conference on Certified Programs and Proofs,
CPP 2019, Cascais, Portugal, January 14-15, 2019. 65--77.
https://doi.org/10.1145/3293880.3294093\\
{[}39{]} C. Kaliszyk and K. Pąk. 2017. Progress in the independent
certification of mizar mathematical library in isabelle. In 2017
Federated Conference on Computer Science and Information Systems
(FedCSIS). 227--236. https://doi.org/10.15439/2017F289\\
{[}40{]} Cezary Kaliszyk, Karol Pąk, and Josef Urban. 2016. Towards a
Mizar Environment for Isabelle: Foundations and Language. In Proceedings
of the 5th ACM SIGPLAN Conference on Certified Programs and Proofs (CPP
2016). ACM, New York, NY, USA, 58--65.
https://doi.org/10.1145/2854065.2854070\\
{[}41{]} Florian Kammüller, Markus Wenzel, and Lawrence C. Paulson.
1999. Locales - A Sectioning Concept for Isabelle. In Theorem Proving in
Higher Order Logics, 12th International Conference, TPHOLs'99, Nice,
France, September, 1999, Proceedings. 149--166.
https://doi.org/10.1007/3-540-48256-3\_11\\
{[}42{]} M. Kaufmann, P. Manolios, and J.S. Moore. 2000. Computer-Aided
Reasoning: ACL2 Case Studies. Advances in Formal Methods, Vol. 4.
Springer US. https://doi.org/10.1007/978-1-4757-3188-0\\
{[}43{]} Holden Lee. 2014. Vector Spaces. Archive of Formal Proofs
(2014). http://isa-afp.org/entries/VectorSpace.html, Formal proof
development.\\
{[}44{]} Robert Y. Lewis. 2019. A formal proof of Hensel's lemma over
the p-adic integers. In Proceedings of the 8th ACM SIGPLAN International
Conference on Certified Programs and Proofs, CPP 2019, Cascais,
Portugal, January 14-15, 2019. 15--26.
https://doi.org/10.1145/3293880.3294089\\
{[}45{]} Paul-Nicolas Madelaine. 2019. Arithmetic and casting in Lean.
Technical Report. Vrije Universiteit Amsterdam.
https://lean-forward.github.io/internships/arithmetic\_and\_casting\_in\_lean.pdf\\
{[}46{]} Assia Mahboubi and Enrico Tassi. 2017. Mathematical Components.
https://math-comp.github.io/mcb/\\
{[}47{]} Norman Megill and David A. Wheeler. 2019. Metamath: A Computer
Language for Mathematical Proofs. Lulu Press.\\
{[}48{]} Walther Neuper. 2019. Technologies for ``Complete, Transparent
\& Interactive Models of Math'' in Education. Electronic Proceedings in
Theoretical Computer Science 290 (03 2019), 76--95.
https://doi.org/10.4204/EPTCS.290.6

{[}49{]} Tobias Nipkow, Lawrence C Paulson, and Markus Wenzel. 2002.
Isabelle/HOL: a proof assistant for higher-order logic. Vol. 2283.
Springer Science \& Business Media.
https://doi.org/10.1007/3-540-45949-9\\
{[}50{]} Sam Owre, John M. Rushby, and Natarajan Shankar. 1992. PVS: A
Prototype Verification System. In Automated Deduction - CADE-11, 11th
International Conference on Automated Deduction, Saratoga Springs, NY,
USA, June 15-18, 1992, Proceedings. 748--752.
https://doi.org/10.1007/3-540-55602-8\_217\\
{[}51{]} Karol Pak. 2016. Topological Foundations for a Formal Theory of
Manifolds. In Joint Proceedings of the FM4M, MathUI, and ThEdu
Workshops, Doctoral Program, and Work in Progress at the Conference on
Intelligent Computer Mathematics 2016 co-located with the 9th Conference
on Intelligent Computer Mathematics (CICM 2016), Bialystok, Poland, July
25-29, 2016. 14--16. http://ceur-ws.org/Vol-1785/F4.pdf\\
{[}52{]} Frank Pfenning and Christine Paulin-Mohring. 1989. Inductively
Defined Types in the Calculus of Constructions. In Mathematical
Foundations of Programming Semantics, 5th International Conference,
Tulane University, New Orleans, Louisiana, USA, March 29 - April 1,
1989, Proceedings. 209--228. https://doi.org/10.1007/BFb0040259\\
{[}53{]} William Pugh. 1991. The Omega Test: A Fast and Practical
Integer Programming Algorithm for Dependence Analysis. In Proceedings of
the 1991 ACM/IEEE Conference on Supercomputing (Supercomputing '91).
ACM, New York, NY, USA, 4--13. https://doi.org/10.1145/125826.125848\\
{[}54{]} Daniel Selsam and Leonardo Moura. 2016. Congruence Closure in
Intensional Type Theory. In Proceedings of the 8th International Joint
Conference on Automated Reasoning - Volume 9706. Springer-Verlag,
Berlin, Heidelberg, 99--115.
https://doi.org/10.1007/978-3-319-40229-1\_8\\
{[}55{]} Konrad Slind and Michael Norrish. 2008. A Brief Overview of
HOL4. In Theorem Proving in Higher Order Logics, Otmane Ait Mohamed,
César Muñoz, and Sofiène Tahar (Eds.). Springer Berlin Heidelberg,
Berlin, Heidelberg, 28--32.
https://doi.org/10.1007/978-3-540-71067-7\_6\\
{[}56{]} Bas Spitters and Eelis van der Weegen. 2011. Type classes for
mathematics in type theory. Mathematical Structures in Computer Science
21, 4 (2011), 795--825. https://doi.org/10.1017/S0960129511000119\\
{[}57{]} Sebastian Ullrich and Leonardo de Moura. 2019. Counting
Immutable Beans: Reference Counting Optimized for Purely Functional
Programming. arXiv:cs.PL/1908.05647\\
{[}58{]} Josef Urban and Geoff Sutcliffe. 2008. ATP-based
Cross-Verification of Mizar Proofs: Method, Systems, and First
Experiments. Mathematics in Computer Science 2, 2 (2008), 231--251.
https://doi.org/10.1007/s11786-008-0053-7\\
{[}59{]} Philip Wadler and Stephen Blott. 1989. How to Make ad-hoc
Polymorphism Less ad-hoc. In Conference Record of the Sixteenth Annual
ACM Symposium on Principles of Programming Languages, Austin, Texas,
USA, January 11-13, 1989. 60--76. https://doi.org/10.1145/75277.75283\\
{[}60{]} Markus Wenzel. 1997. Type Classes and Overloading in
Higher-Order Logic. In Theorem Proving in Higher Order Logics, 10th
International Conference, TPHOLs'97, Murray Hill, NJ, USA, August 19-22,
1997, Proceedings. 307--322. https://doi.org/10.1007/BFb0028402\\
{[}61{]} H. P. Williams. 1986. Fourier's Method of Linear Programming
and Its Dual. The American Mathematical Monthly 93, 9 (1986), 681--695.
http://www.jstor.org/stable/2322281\\
{[}62{]} Minchao Wu and Rajeev Goré. 2019. Verified Decision Procedures
for Modal Logics. In 10th International Conference on Interactive
Theorem Proving (ITP 2019) (Leibniz International Proceedings in
Informatics (LIPIcs)), John Harrison, John O'Leary, and Andrew Tolmach
(Eds.), Vol. 141. Schloss Dagstuhl--Leibniz-Zentrum fuer Informatik,
Dagstuhl, Germany, 31:1--31:19.
https://doi.org/10.4230/LIPIcs.ITP.2019.31

\end{document}
